\chapter{2020年上海卷（春）}

\section{填空题}

\begin{problem}
集合 \(A=\{1,3\}\)，\(B=\{1,2,a\}\)，若 \(A\subseteq B\)，则 \(a=\fillinblank{}\).
\end{problem}

\begin{answer}
\(3\)
\end{answer}

\begin{solution}
因为 \(3\in A\)，且 \(A\subseteq B\)，所以 \(3\in B\)，故 \(a=3\).
\end{solution}

\begin{problem}
不等式 \(\dfrac{1}{x}>3\) 的解集为 \(\fillinblank{}\).
\end{problem}

\begin{answer}
\(\left(0,\dfrac{1}{3}\right)\)
\end{answer}

\begin{solution}
由 \(\dfrac{1}{x}>3\) 得 \(\dfrac{1-3x}{x}>0\)，则 \(x(1-3x)>0\)，即 \(x(3x-1)<0\)，解得 \(0<x<\dfrac{1}{3}\)，所以不等式的解集是 \(\left(0,\dfrac{1}{3}\right)\).
\end{solution}

\begin{problem}
函数 \(y=\tan 2x\) 的最小正周期为 \(\fillinblank{}\).
\end{problem}

\begin{answer}
\(\dfrac{\pi}{2}\)
\end{answer}

\begin{solution}
函数 \(y=\tan 2x\) 的最小正周期为 \(\dfrac{\pi}{2}\).
\end{solution}

\begin{problem}
已知复数 \(z\) 满足 \(z+2\bar z=6+\i\)，则 \(z\) 的实部为 \(\fillinblank{}\).
\end{problem}

\begin{answer}
\(2\)
\end{answer}

\begin{solution}
设 \(z=a+b\i\)，（\(a\)，\(b\in\R\)）. 因为复数 \(z\) 满足 \(z+2\bar z=6+\i\)，所以 \(3a-b\i=6+\i\)，可得 \(3a=6\)，\(-b=1\)，解得 \(a=2\)，\(b=-1\). 则 \(z\) 的实部为 \(2\).
\end{solution}

\begin{problem}
已知 \(3\sin 2x=2\sin x\)，\(x\in(0,\pi)\)，则 \(x=\fillinblank{}\).
\end{problem}

\begin{answer}
\(\arccos\dfrac{1}{3}\)
\end{answer}

\begin{solution}
因为 \(3\sin 2x=2\sin x\)，所以 \(6\sin x\cos x=2\sin x\). 因为 \(x\in(0,\pi)\)，所以 \(\sin x\ne0\)，从而 \(\cos x=\dfrac{1}{3}\)，故 \(x=\arccos\dfrac{1}{3}\).
\end{solution}

\begin{problem}
若函数 \(y=a\cdot3^x+\dfrac{1}{3^x}\) 为偶函数，则 \(a=\fillinblank{}\).
\end{problem}

\begin{answer}
\(1\)
\end{answer}

\begin{solution}
根据题意，函数 \(f(x)=a\cdot3^x+\dfrac{1}{3^x}\) 为偶函数，则 \(f(-x)=f(x)\)，即 \(a\cdot3^{-x}+\dfrac{1}{3^{-x}}=a\cdot3^x+\dfrac{1}{3^x}\). 变形可得 \(a(3^x-3^{-x})=3^x-3^{-x}\)，必有 \(a=1\).
\end{solution}

\begin{problem}
已知直线 \(l_1:x+ay=1\)，\(l_2:ax+y=1\)，若 \(l_1\parallel l_2\)，则 \(l_1\) 与 \(l_2\) 的距离为 \(\fillinblank{}\).
\end{problem}

\begin{answer}
\(\sqrt{2}\)
\end{answer}

\begin{solution}
直线 \(l_1:x+ay=1\)，\(l_2:ax+y=1\). 当 \(l_1\parallel l_2\) 时，\(a^2-1=0\)，解得 \(a=\pm1\). 当 \(a=1\) 时 \(l_1\) 与 \(l_2\) 重合，不满足题意；当 \(a=-1\) 时 \(l_1:x-y-1=0\)，\(l_2:x-y+1=0\)，则 \(l_1\) 与 \(l_2\) 的距离为 \(\dfrac{|-1-1|}{\sqrt{1^2+(-1)^2}}=\sqrt{2}\).
\end{solution}

\begin{problem}
已知二项式 \((2x+\sqrt{x})^5\)，则展开式中 \(x^3\) 的系数为 \(\fillinblank{}\).
\end{problem}

\begin{answer}
\(10\)
\end{answer}

\begin{solution}
由二项式展开式的通项，\(T_{r+1}=\mathrm C_5^r(2x)^{5-r}(\sqrt{x})^r=\mathrm C_5^r2^{5-r}x^{5-\frac{r}{2}}\). 令 \(5-\dfrac{r}{2}=3\)，得 \(r=4\)，代入通项得 \(\mathrm C_5^4 2x^3=10x^3\)，所以展开式中 \(x^3\) 的系数为 \(10\).
\end{solution}

\begin{problem}
\(\triangle ABC\) 中，\(D\) 是 \(BC\) 中点，\(AB=2\)，\(BC=3\)，\(AC=4\)，则 \(\overrightarrow{AD}\cdot\overrightarrow{AB}=\fillinblank{}\).
\end{problem}

\begin{answer}
\(\dfrac{19}{4}\)
\end{answer}

\begin{solution}
在 \(\triangle ABC\) 中，\(AB=2\)，\(BC=3\)，\(AC=4\). 由余弦定理得 \(\cos\angle BAC=\dfrac{AB^2+AC^2-BC^2}{2AB\cdot AC}=\dfrac{4+16-9}{2\times2\times4}=\dfrac{11}{16}\)，所以 \(\overrightarrow{AB}\cdot\overrightarrow{AC}=2\times4\times\dfrac{11}{16}=\dfrac{11}{2}\). 又 \(D\) 是 \(BC\) 的中点，所以 \(\overrightarrow{AD}=\dfrac{1}{2}(\overrightarrow{AB}+\overrightarrow{AC})\). 因此 \(\overrightarrow{AD}\cdot\overrightarrow{AB}=\dfrac{1}{2}(\overrightarrow{AB}^{\,2}+\overrightarrow{AB}\cdot\overrightarrow{AC})=\dfrac{1}{2}\left(4+\dfrac{11}{2}\right)=\dfrac{19}{4}\).
\end{solution}

\begin{problem}
已知 \(A=\{-3,-2,-1,0,1,2,3\}\)，\(a\)，\(b\in A\)，则 \(|a|<|b|\) 的情况有 \(\fillinblank{}\) 种.
\end{problem}

\begin{answer}
\(18\)
\end{answer}

\begin{solution}
当 \(a=-3\) 时，有 \(0\) 种；当 \(a=-2\) 时，有 \(2\) 种；当 \(a=-1\) 时，有 \(4\) 种；当 \(a=0\) 时，有 \(6\) 种；当 \(a=1\) 时，有 \(4\) 种；当 \(a=2\) 时，有 \(2\) 种；当 \(a=3\) 时，有 \(0\) 种，故共有 \(2+4+6+4+2=18\) 种.
\end{solution}

\begin{problem}
已知 \(A_1\)，\(A_2\)，\(A_3\)，\(A_4\)，\(A_5\) 五个点，满足 \(\overrightarrow{A_nA_{n+1}}\cdot\overrightarrow{A_{n+1}A_{n+2}}=0\) （\(n=1\)，\(2\)，\(3\)），\(|A_nA_{n+1}|\cdot|A_{n+1}A_{n+2}|=n+1\) （\(n=1\)，\(2\)，\(3\)），则 \(|A_1A_5|\) 的最小值为 \(\fillinblank{}\).
\end{problem}

\begin{answer}
\(\dfrac{\sqrt{6}}{3}\)
\end{answer}

\begin{solution}
设 \(|A_1A_2|=x\)，则 \(|A_2A_3|=\dfrac{2}{x}\)，\(|A_3A_4|=\dfrac{3x}{2}\)，\(|A_4A_5|=\dfrac{8}{3x}\). 设 \(A_1(0,0)\)，如图.

\solutionfigure{\bitmapfigure[max width=0.50\linewidth]{img/2020/shanghai_spring/q11_solution_fig1.png}}

因为求 \(|A_1A_5|\) 的最小值，所以可设 \(A_2(x,0)\)，\(A_3\left(x,\dfrac{2}{x}\right)\)，\(A_4\left(-\dfrac{x}{2},\dfrac{2}{x}\right)\)，\(A_5\left(-\dfrac{x}{2},-\dfrac{2}{3x}\right)\). 因此 \(|A_1A_5|^2=\left(-\dfrac{x}{2}\right)^2+\left(-\dfrac{2}{3x}\right)^2=\dfrac{x^2}{4}+\dfrac{4}{9x^2}\geq\dfrac{\sqrt{6}}{3}\,^2\)，当且仅当 \(\dfrac{x^2}{4}=\dfrac{4}{9x^2}\)，即 \(x=\dfrac{2\sqrt{3}}{3}\) 时取等号，故 \(|A_1A_5|\) 的最小值为 \(\dfrac{\sqrt{6}}{3}\).
\end{solution}

\begin{problem}
已知 \(f(x)=\sqrt{x-1}\)，其反函数为 \(f^{-1}(x)\)，若 \(f^{-1}(x)-a=f(x+a)\) 有实数根，则 \(a\) 的取值范围为 \(\fillinblank{}\).
\end{problem}

\begin{answer}
\(\left[\dfrac{3}{4},+\infty\right)\)
\end{answer}

\begin{solution}
因为 \(y=f^{-1}(x)-a\) 与 \(y=f(x+a)\) 互为反函数，又方程 \(f^{-1}(x)-a=f(x+a)\) 有实数根，则函数 \(y=f(x+a)\) 的图象与直线 \(y=x\) 有交点. 所以 \(\sqrt{x+a-1}=x\) 有根，即 \(a=x^2-x+1=\left(x-\dfrac{1}{2}\right)^2+\dfrac{3}{4}\)，故答案为 \(\left[\dfrac{3}{4},+\infty\right)\).

\solutionfigure{\bitmapfigure[max width=0.50\linewidth]{img/2020/shanghai_spring/q12_fig1.png}}
\end{solution}

\section{单选题}

\begin{problem}
计算 \(\displaystyle\lim_{n\to\infty}\dfrac{3^n+5^n}{3^{n-1}+5^{n-1}}=\)（\quad）.

\choices{\(3\)}{\(\dfrac{5}{3}\)}{\(\dfrac{3}{5}\)}{\(5\)}
\end{problem}

\begin{answer}
D
\end{answer}

\begin{solution}
分子分母同时除以 \(5^{n-1}\)，则 \(\displaystyle\lim_{n\to\infty}\dfrac{3^n+5^n}{3^{n-1}+5^{n-1}}=\lim_{n\to\infty}\dfrac{3\left(\dfrac{3}{5}\right)^{n-1}+5}{\left(\dfrac{3}{5}\right)^{n-1}+1}=5\)，故选 D.
\end{solution}

\begin{problem}
“\(\alpha=\beta\)”是“\(\sin^2\alpha+\cos^2\beta=1\)”的（\quad）.

\choices{充分非必要条件}{必要非充分条件}{充要条件}{既非充分又非必要条件}
\end{problem}

\begin{answer}
A
\end{answer}

\begin{solution}
若 \(\alpha=\beta\)，则 \(\sin^2\alpha+\cos^2\beta=\sin^2\alpha+\cos^2\alpha=1\)，所以“\(\alpha=\beta\)”是“\(\sin^2\alpha+\cos^2\beta=1\)”的充分条件. 若 \(\sin^2\alpha+\cos^2\beta=1\)，则 \(\sin^2\alpha=\sin^2\beta\)，不能推出 \(\alpha=\beta\)，所以“\(\alpha=\beta\)”不是“\(\sin^2\alpha+\cos^2\beta=1\)”的必要条件.
故“\(\alpha=\beta\)”是“\(\sin^2\alpha+\cos^2\beta=1\)”的充分非必要条件，故选 A.
\end{solution}

\begin{problem}
已知椭圆 \(\dfrac{x^2}{2}+y^2=1\)，作垂直于 \(x\) 轴的垂线交椭圆于 \(A\)，\(B\) 两点，作垂直于 \(y\) 轴的垂线交椭圆于 \(C\)，\(D\) 两点，且 \(AB=CD\)，两垂线相交于点 \(P\)，则点 \(P\) 的轨迹是（\quad）.

\choices{椭圆}{双曲线}{圆}{抛物线}
\end{problem}

\begin{answer}
B
\end{answer}

\begin{solution}
由 \(AB=2\sqrt{1-\dfrac{m^2}{2}}\leq2\)，且 \(AB=CD\)，可知两条弦长相等；设 \(A(m,t)\)，\(D(t,n)\)，则 \(P(m,n)\). 因为 \(\dfrac{m^2}{2}+t^2=1\)，\(\dfrac{t^2}{2}+n^2=1\)，消去 \(t\) 可得 \(2n^2-\dfrac{m^2}{2}=1\)，故选 B.

\FigureLayoutDeclare{img/2020/shanghai_spring/q15_fig1.png}{5.0cm}{32bp 0bp 0bp 0bp}
\solutionfigure{\bitmapfigure[max width=0.50\linewidth]{img/2020/shanghai_spring/q15_fig1.png}}
\end{solution}

\begin{problem}
数列 \(\{a_n\}\) 各项均为实数，对任意 \(n\in\mathbb N^*\) 满足 \(a_{n+3}=a_n\)，且行列式 \(\begin{vmatrix}a_n&a_{n+1}\\a_{n+2}&a_{n+3}\end{vmatrix}=c\) 为定值，则下列选项中不可能的是（\quad）.

\choices{\(a_1=1,c=1\)}{\(a_1=2,c=2\)}{\(a_1=-1,c=4\)}{\(a_1=2,c=0\)}
\end{problem}

\begin{answer}
B
\end{answer}

\begin{solution}
行列式 \(\begin{vmatrix}a_n&a_{n+1}\\a_{n+2}&a_{n+3}\end{vmatrix}=a_na_{n+3}-a_{n+1}a_{n+2}=c\). 由于对任意 \(n\in\mathbb N^*\) 满足 \(a_{n+3}=a_n\)，所以
\[
\begin{cases}
a_n^2-a_{n+1}a_{n+2}=c,\\
a_{n+1}^2-a_{n+2}a_{n+3}=c.
\end{cases}
\]
作差整理得：\(a_{n+1}=a_n\)（常数列，\(c=0\)），或 \(a_n+a_{n+1}+a_{n+2}=0\). 当 \(a_n+a_{n+1}+a_{n+2}=0\) 时，\(a_{n+1}+a_{n+2}=-a_n\) 且 \(a_{n+1}a_{n+2}=a_n^2-c\)，所以方程 \(X^2+a_nX+a_n^2-c=0\) 有两根 \(a_{n+1}\)，\(a_{n+2}\)，从而 \(\Delta=a_n^2-4(a_n^2-c)=4c-3a_n^2\geqslant0\). 等号情形也可能出现，但选项 B 仍不可能，故选 B.
\end{solution}

\section{解答题}

\begin{problem}
已知四棱锥 \(P-ABCD\)，底面 \(ABCD\) 为正方形，边长为 \(3\)，\(PD\perp\) 平面 \(ABCD\).
\begin{enumerate}
\item 若 \(PC=5\)，求四棱锥 \(P-ABCD\) 的体积；
\item 若直线 \(AD\) 与 \(BP\) 的夹角为 \(60^\circ\)，求 \(PD\) 的长.
\end{enumerate}

\FigureLayoutDeclare{img/2020/shanghai_spring/q17_fig1.png}{4.5cm}{13bp 0bp 25bp 0bp}
\bitmapfigure[max width=0.50\linewidth]{img/2020/shanghai_spring/q17_fig1.png}
\end{problem}

\begin{solution}
\begin{enumerate}
\item 因为 \(PD\perp\) 平面 \(ABCD\)，所以 \(PD\perp DC\). 因为 \(CD=3\)，\(PC=5\)，所以 \(PD=4\)，故 \(V_{P-ABCD}=\dfrac{1}{3}\times3^2\times4=12\). 所以四棱锥 \(P-ABCD\) 的体积为 \(12\).
\item 因为 \(ABCD\) 是正方形，\(PD\perp\) 平面 \(ABCD\)，所以 \(BC\perp PD\)，\(BC\perp CD\). 又 \(PD\cap CD=D\)，所以 \(BC\perp\) 平面 \(PCD\)，从而 \(BC\perp PC\). 因为异面直线 \(AD\) 与 \(PB\) 所成角为 \(60^\circ\)，且 \(BC\parallel AD\)，所以 \(\angle PBC=60^\circ\). 在 \(\mathrm{Rt}\triangle PBC\) 中，\(BC=3\)，故 \(PC=3\sqrt{3}\). 在 \(\mathrm{Rt}\triangle PDC\) 中，又 \(CD=3\)，所以 \(PD=3\sqrt{2}\).
\end{enumerate}
\end{solution}

\begin{problem}
已知各项均为正数的数列 \(\{a_n\}\)，其前 \(n\) 项和为 \(S_n\)，\(a_1=1\).
\begin{enumerate}
\item 若数列 \(\{a_n\}\) 为等差数列，\(S_{10}=70\)，求数列 \(\{a_n\}\) 的通项公式；
\item 若数列 \(\{a_n\}\) 为等比数列，\(a_4=\dfrac{1}{8}\)，求满足 \(S_n>100a_n\) 时 \(n\) 的最小值.
\end{enumerate}
\end{problem}

\begin{solution}
\begin{enumerate}
\item 设等差数列的公差为 \(d\). 由 \(S_{10}=70\)，得 \(10+\dfrac{1}{2}\times10\times9d=70\)，解得 \(d=\dfrac{4}{3}\)，则 \(a_n=1+\dfrac{4}{3}(n-1)=\dfrac{4}{3}n-\dfrac{1}{3}\).
\item 设等比数列的公比为 \(q\). 由 \(a_4=\dfrac{1}{8}\)，\(a_1=1\)，可得 \(q^3=\dfrac{1}{8}\)，即 \(q=\dfrac{1}{2}\). 则 \(a_n=\left(\dfrac{1}{2}\right)^{n-1}\)，\(S_n=\dfrac{1-\left(\frac{1}{2}\right)^n}{1-\frac{1}{2}}=2-\left(\dfrac{1}{2}\right)^{n-1}\). \(S_n>100a_n\)，即 \(2-\left(\dfrac{1}{2}\right)^{n-1}>100\cdot\left(\dfrac{1}{2}\right)^{n-1}\)，即 \(2^n>101\)，可得 \(n\geq7\)，即 \(n\) 的最小值为 \(7\).
\end{enumerate}
\end{solution}

\begin{problem}
有一条长为 \(120\) 米的步行道 \(OA\)，\(A\) 是垃圾投放点 \(\omega_1\)，若以 \(O\) 为原点，\(OA\) 为 \(x\) 轴正半轴建立直角坐标系，设点 \(B(x,0)\)，现要建设另一座垃圾投放点 \(\omega_2(t,0)\)，函数 \(f_t(x)\) 表示与 \(B\) 点距离最近的垃圾投放点的距离.
\begin{enumerate}
\item 若 \(t=60\)，求 \(f_{60}(10)\)，\(f_{60}(80)\)，\(f_{60}(95)\) 的值，并写出 \(f_{60}(x)\) 的函数解析式；
\item 若可以通过 \(f_t(x)\) 与坐标轴围成的面积来测算扔垃圾的便利程度，面积越小越便利. 问：垃圾投放点 \(\omega_2\) 建在何处才能比建在中点时更加便利？
\end{enumerate}
\end{problem}

\begin{solution}
\begin{enumerate}
\item 投放点 \(\omega_1(120,0)\)，\(\omega_2(60,0)\). 由题意得 \(f_{60}(x)=\min\{|60-x|,|120-x|\}\)，所以 \(f_{60}(10)=50\)，\(f_{60}(80)=20\)，\(f_{60}(95)=25\). 当 \(x\leq90\) 时，\(f_{60}(x)=|60-x|\)；当 \(x>90\) 时，\(f_{60}(x)=|120-x|\). 综上，
\[
f_{60}(x)=
\begin{cases}
|60-x|,&x\leq90,\\
|120-x|,&x>90.
\end{cases}
\]
\item 由题意得 \(f_t(x)=\min\{|t-x|,|120-x|\}\)，所以
\[
f_t(x)=
\begin{cases}
|t-x|,&x\leq0.5(120+t),\\
|120-x|,&x>0.5(120+t).
\end{cases}
\]
因此 \(f_t(x)\) 与坐标轴围成的面积为 \(S=\dfrac{1}{2}t^2+\dfrac{1}{4}(120-t)^2=\dfrac{3}{4}t^2-60t+3600\). 由题意，\(S<S(60)\)，即 \(\dfrac{3}{4}t^2-60t+3600<2700\)，解得 \(20<t<60\)，即垃圾投放点 \(\omega_2\) 建在 \((20,0)\) 与 \((60,0)\) 之间时，比建在中点时更加便利.
\end{enumerate}
\solutionfigure{\bitmapfigure[max width=0.50\linewidth]{img/2020/shanghai_spring/q19_solution_fig1.png}}
\end{solution}

\begin{problem}
已知抛物线 \(y^2=x\) 上的动点 \(M(x_0,y_0)\)，过 \(M\) 分别作两条直线交抛物线于 \(P\)，\(Q\) 两点，交直线 \(x=t\) 于 \(A\)，\(B\) 两点.
\begin{enumerate}
\item 若点 \(M\) 纵坐标为 \(\sqrt{2}\)，求 \(M\) 与焦点的距离；
\item 若 \(t=-1\)，\(P(1,1)\)，\(Q(1,-1)\)，求证：\(y_A\cdot y_B\) 为常数；
\item 是否存在 \(t\)，使得 \(y_A\cdot y_B=1\) 且 \(y_P\cdot y_Q\) 为常数？若存在，求出 \(t\) 的所有可能值；若不存在，请说明理由.
\end{enumerate}
\end{problem}

\begin{solution}
\begin{enumerate}
\item 抛物线 \(y^2=x\) 上的动点为 \(M(x_0,y_0)\)，且点 \(M\) 纵坐标为 \(\sqrt{2}\)，所以点 \(M\) 的横坐标 \(x_M=(\sqrt{2})^2=2\). 由 \(y^2=x\)，得 \(p=\dfrac{1}{2}\)，所以 \(M\) 与焦点的距离为 \(MF=x_M+\dfrac{p}{2}=2+\dfrac{1}{4}=\dfrac{9}{4}\).
\item 设 \(M(y_0^2,y_0)\)，直线 \(PM\colon y-1=\dfrac{y_0-1}{y_0^2-1}(x-1)\). 当 \(x=-1\) 时，\(y_A=\dfrac{y_0-1}{y_0+1}\). 直线 \(QM\colon y+1=\dfrac{y_0+1}{y_0^2-1}(x-1)\)，当 \(x=-1\) 时，\(y_B=\dfrac{-y_0-1}{y_0-1}\)，所以 \(y_Ay_B=-1\)，故 \(y_A\cdot y_B\) 为常数 \(-1\).
\item 设 \(M(y_0^2,y_0)\)，\(A(t,y_A)\)，直线 \(MA\colon y-y_0=\dfrac{y_0-y_A}{y_0^2-t}(x-y_0^2)\). 联立 \(y^2=x\)，得
\[
y^2-\dfrac{y_0^2-t}{y_0-y_A}y+\dfrac{y_0^2-t}{y_0-y_A}y_0-y_0^2=0
\]
所以 \(y_0+y_P=\dfrac{y_0^2-t}{y_0-y_A}\)，即 \(y_P=\dfrac{y_0y_A-t}{y_0-y_A}\). 同理得 \(y_Q=\dfrac{y_0y_B-t}{y_0-y_B}\). 因为 \(y_A\cdot y_B=1\)，所以
\[
y_Py_Q=\dfrac{y_0^2-ty_0(y_A+y_B)+t^2}{y_0^2-y_0(y_A+y_B)+1}
\]
要使 \(y_Py_Q\) 为常数，即 \(t=1\)，此时 \(y_Py_Q\) 为常数 \(1\). 所以存在 \(t=1\)，使得 \(y_A\cdot y_B=1\) 且 \(y_P\cdot y_Q\) 为常数 \(1\).
\end{enumerate}
\end{solution}

\begin{problem}
已知非空集合 \(A\subseteq\R\)，函数 \(y=f(x)\) 的定义域为 \(D\)，若对任意 \(t\in A\) 且 \(x\in D\)，不等式 \(f(x)\leq f(x+t)\) 恒成立，则称函数 \(f(x)\) 具有 \(A\) 性质.
\begin{enumerate}
\item 当 \(A=\{-1\}\)，判断 \(f(x)=-x\)，\(g(x)=2^x\) 是否具有 \(A\) 性质；
\item 当 \(A=(0,1)\)，\(f(x)=x+\dfrac{1}{x}\)，\(x\in[a,+\infty)\)，若 \(f(x)\) 具有 \(A\) 性质，求 \(a\) 的取值范围；
\item 当 \(A=\{-2,m\}\)，\(m\in\Z\)，若 \(D\) 为整数集且具有 \(A\) 性质的函数均为常值函数，求所有符合条件的 \(m\) 的值.
\end{enumerate}
\end{problem}

\begin{solution}
\begin{enumerate}
\item 因为 \(f(x)=-x\) 为减函数，所以 \(f(x)<f(x-1)\)，故 \(f(x)=-x\) 具有 \(A\) 性质；\(g(x)=2^x\) 为增函数，所以 \(g(x)>g(x-1)\)，故 \(g(x)=2^x\) 不具有 \(A\) 性质.
\item 依题意，对任意 \(t\in(0,1)\)，应有 \(f(x)\leq f(x+t)\)（\(x\geq a\)）. 函数 \(f(x)=x+\dfrac1x\) 在 \((0,1)\) 上递减、在 \((1,+\infty)\) 上递增，因此当且仅当 \(a\geq1\) 时，函数在 \([a,+\infty)\) 上单调递增，满足 A 性质. 故 \(a\in[1,+\infty)\).
\item 由 \(t=-2\) 得 \(f(x)\leq f(x-2)\)，而题设说明具有 \(A\) 性质的函数均为常值函数，因此 \(f(x)=f(x-2)\) 恒成立，函数以 \(2\) 为周期. 设 \(f(2k)=p\)，\(f(2n-1)=q\)（\(n,k\in\Z\)），则需由 \(m\) 使两个奇偶类相互联系；这等价于 \(m\) 为奇数. 反之，任意奇数 \(m\) 都能使步长 \(-2\) 与 \(m\) 生成所有整数位置，故所有满足条件的 \(m\) 为奇数.
\end{enumerate}
\end{solution}
